\subsection{Calculus Vol2 - 4.1 Basics of Differential Equations}

\url{https://openstax.org/books/calculus-volume-2/pages/4-1-basics-of-differential-equations}

\begin{enumerate}
  \item What is a differential equation?
  \item What is meant by a solution to a differential equation?
  \item How can you verify that a function $y=f(x)$ is a solution to a given differential equation?
  \item Give an example of a differential equation and one of its solutions.
  \item Why are solutions to differential equations often not unique?
  \item What property of derivatives explains the non-uniqueness of solutions?
  \item What is the order of a differential equation?
  \item How do you determine the order of a differential equation from its expression?
  \item What is the order of the differential equation $y' - 4y = x^2 - 3x + 4$?
  \item What is the order of the differential equation $x^2 y''' - 3x y'' + x y' - 3y = \sin x$?
  \item What is a general solution of a differential equation?
  \item Why does a general solution typically include an arbitrary constant $C$?
  \item For the differential equation $y' = 2x$, what is its general solution?
  \item What does the family of functions $y = x^2 + C$ represent?
  \item What is a particular solution of a differential equation?
  \item How do you obtain a particular solution from a general solution?
  \item What is an initial-value problem?
  \item How does an initial condition help determine a unique solution?
  \item Solve the initial-value problem:
    \[
      y' = 2x, \quad y(2) = 7.
    \]
  \item What is the difference between a general solution and a particular solution?
  \item What does it mean for a function to satisfy a differential equation?
  \item Verify that $y = e^{-3x} + 2x + 3$ is a solution of
    \[
      y' + 3y = 6x + 11.
    \]
  \item Verify that $y = 2e^{3x} - 2x - 2$ is a solution of
    \[
      y' - 3y = 6x + 4.
    \]
  \item When checking a solution, why must both $y$ and its derivatives be substituted into the equation?
  \item Can two different functions be solutions to the same differential equation? Explain why.
  \item What role does the constant $C$ play geometrically in the family of solutions?
  \item How would the graph of $y = x^2 + C$ change as $C$ varies?
  \item Why is identifying the order of a differential equation useful?
  \item What information is needed in addition to a differential equation to determine a unique solution?
\end{enumerate}
